% from ega2-1-fr.tex
\oldpage[II]{165}
Lorsque $Y=\operatorname{Spec}(A)$ est affine, on a, avec les notations de
(\hyperref[II.8.3.2-fr]{8.3.2})
\begin{equation}
\label{II.8.3.4.4-fr}
  E=\bigcup C_f,\qquad
  \widehat{E}=\bigcup\widehat{C}_f,\qquad
  C_f=C\cap\widehat{C}_f
\tag{8.3.4.4}
\end{equation}
$f$ parcourant l'ensemble des éléments homogènes de $S_+$ (ou seulement une
partie $M$ de cet ensemble engendrant un idéal de $S_+$ dont la racine dans
$S_+$ est $S_+$ lui-même, autrement dit, telle que les $X_f$ pour $f\in M$
recouvrent $X$ (\hyperref[II.2.3.14-fr]{2.3.14})). Le recollement de $C$ et de
$\widehat{C}_f$ le long de $C_f$ est alors déterminé par des morphismes
d'injection $C_f\to C$, $C_f\to\widehat{C}_f$, qui, ainsi qu'on l'a vu
(\hyperref[II.8.3.2.6-fr]{8.3.2.6}) correspondent respectivement aux
homomorphismes canoniques $S\to S_f$, $S_f^{\leq}\to S_f$.
\end{env}

\begin{proposition}[8.3.5]
\label{II.8.3.5-fr}
Avec les notations de (\hyperref[II.8.3.1-fr]{8.3.1}) et
(\hyperref[II.8.3.4-fr]{8.3.4}), le morphisme associé
(\hyperref[II.3.5.1-fr]{3.5.1}) à l'injection canonique
$\varphi:\mathcal{S}\to\widehat{\mathcal{S}}=\mathcal{S}[z]$ est un morphisme
affine surjectif (dit \emph{rétraction canonique})
\begin{equation}
\label{II.8.3.5.1-fr}
  p:\widehat{E}\to X
\tag{8.3.5.1}
\end{equation}
tel que l'on ait
\begin{equation}
\label{II.8.3.5.2-fr}
  p\circ j=1_X.
\tag{8.3.5.2}
\end{equation}
\end{proposition}

Pour démontrer la proposition, on peut se borner au cas où $Y$ est affine.
Compte tenu de l'expression (\hyperref[II.8.3.4.4-fr]{8.3.4.4}) de
$\widehat{E}$, le fait que le domaine de définition $G(\varphi)$ de $p$ est
égal à $\widehat{E}$ résultera de la première des assertions suivantes~:

\begin{env}[8.3.5.3]
\label{II.8.3.5.3-fr}
Si $Y=\operatorname{Spec}(A)$ est affine, et $\mathcal{S}=\widetilde{S}$, on
a, pour tout $f\in S_+$ homogène
\begin{equation}
\label{II.8.3.5.4-fr}
  p^{-1}(X_f)=\widehat{C}_f
\tag{8.3.5.4}
\end{equation}
et la restriction de $p$ à
$\widehat{C}_f=\operatorname{Spec}(S_f^{\leq})$, considérée comme morphisme de
$\widehat{C}_f$ dans $X_f$, correspond à l'injection canonique
$S_{(f)}\to S_f^{\leq}$. Si de plus $f\in S_1$, $\widehat{C}_f$ est isomorphe
à $X_f\otimes_{\mathbf{Z}}\mathbf{Z}[T]$ ($T$ indéterminée).
\end{env}

En effet, la formule (\hyperref[II.8.3.5.4-fr]{8.3.5.4}) n'est qu'un cas
particulier de (\hyperref[II.2.8.1.1-fr]{2.8.1.1}) et la seconde assertion
n'est autre que la définition de $\operatorname{Proj}(\varphi)$ lorsque $Y$
est affine (\hyperref[II.2.8.1-fr]{2.8.1}). Alors la formule
(\hyperref[II.8.3.5.2-fr]{8.3.5.2}) et le fait que $p$ est surjectif résultent
de ce que le composé $\mathcal{S}\to\widehat{\mathcal{S}}\to\mathcal{S}$ des
homomorphismes canoniques est l'identité dans $\mathcal{S}$. Enfin, la
dernière assertion de (\hyperref[II.8.3.5.3-fr]{8.3.5.3}) provient de ce que
$S_f^{\leq}$ est isomorphe à $S_{(f)}[T]$ lorsque $f\in S_1$
(\hyperref[II.2.2.1-fr]{2.2.1}).

\begin{corollary}[8.3.6]
\label{II.8.3.6-fr}
La restriction
\begin{equation}
\label{II.8.3.6.1-fr}
  \pi:E\to X
\tag{8.3.6.1}
\end{equation}
de $p$ à $E$ est un morphisme affine surjectif. Si $Y$ est affine et $f$
homogène dans $S_+$, on a
\begin{equation}
\label{II.8.3.6.2-fr}
  \pi^{-1}(X_f)=C_f
\tag{8.3.6.2}
\end{equation}
et la restriction de $\pi$ à $C_f$ correspond à l'injection canonique
$S_{(f)}\to S_f$. Si en outre $f\in S_1$, $C_f$ est isomorphe à
$X_f\otimes_{\mathbf{Z}}\mathbf{Z}[T,T^{-1}]$ ($T$ indéterminée).
\end{corollary}

La formule (\hyperref[II.8.3.6.2-fr]{8.3.6.2}) résulte aussitôt de
(\hyperref[II.8.3.5.3-fr]{8.3.5.3}) et
(\hyperref[II.8.3.2.5-fr]{8.3.2.5}), et démontre la surjectivité de $\pi$~; on
a vu par ailleurs que l'immersion $i$, restreinte à $C_f$, correspondait

